% =========================================================
% INTRODUCCIÓN — versión RUCP (~1,500 palabras)
% =========================================================

\section{Introducción}
\label{sec:introduccion}

El 24 de noviembre de 2024, Yamandú Orsi (Frente Amplio, FA) derrotó a Álvaro Delgado (Partido Nacional, PN) en el balotaje presidencial uruguayo por aproximadamente 95{,}000 votos, poniendo fin a cinco años de gobierno de la Coalición Republicana. El resultado sorprendió a analistas que estimaban que la coalición ---integrada por el PN, el Partido Colorado (PC), Cabildo Abierto (CA) y el Partido Independiente (PI)--- capitalizaría su organización en segunda vuelta. Un análisis posterior revela un patrón inesperado que motiva este estudio: los votantes rurales de CA transfirieron su apoyo al FA en proporciones \textit{mayores} que los votantes urbanos del mismo partido, contradiciendo la expectativa convencional de mayor conservadurismo en zonas rurales para un partido de perfil militar-populista. Se denomina este fenómeno ``rebelión rural'' y constituye la tesis central del trabajo.

La evidencia de este patrón no es accesible de forma directa. El secreto del voto impide conocer las decisiones individuales; solo se dispone de recuentos de votos agregados a nivel de circuito electoral. El desafío de inferir comportamiento individual a partir de datos agregados ---el \textit{problema de la inferencia ecológica} \citep{king1997}--- es la razón por la que la \textit{falacia ecológica} \citep{robinson1950ecological} acecha a los análisis de transferencias que utilizan correlaciones agregadas sin corrección estadística. El trabajo fundacional de \citet{king1997} propuso un marco bayesiano para resolver este problema en el caso binario; \citet{rosen2001bayesian} extendieron el enfoque al caso multidimensional $R \times C$ requerido por sistemas multipartidistas como el uruguayo. Trabajo reciente de \citet{greiner2009rxc} evaluó los supuestos de identificación de estos modelos, e \citet{imai2008bayesian} propusieron formulaciones alternativas. En América Latina, las aplicaciones sistemáticas de estos métodos a nivel circuital siguen siendo escasas \citep{calvo2003local}, confiriéndole al presente estudio un carácter exploratorio regionalmente relevante.

El caso de CA merece atención especial en la literatura. Fundado en 2019 por el general retirado Guido Manini Ríos, irrumpió en el sistema de partidos con el 11\% de los votos y colapsó al 2.7\% en 2024 ($-$78\% en caudal electoral). La comprensión del destino de sus votantes remanentes en el balotaje ---y de por qué los circuitos rurales exhiben patrones distintos--- resulta decisiva para interpretar la estabilidad del sistema de partidos uruguayo, uno de los más institucionalizados de América Latina \citep{mainwaring1995party, sartori1976parties}. La consolidación del FA como actor central del sistema ha sido documentada por \citet{perezbentancur2019frente}; el rol de las alianzas coalicionales en el presidencialismo latinoamericano ha sido analizado por \citet{chasquetti2004democracia}.

La literatura sobre divisiones urbano-rurales sugiere que los votantes rurales tienden a preferir partidos de derecha, particularmente en contextos donde la economía agrícola es central \citep{rodden2019s, gonzalez1991electoral}. La segunda vuelta, sin embargo, introduce dinámicas propias: la elección binaria fuerza un reposicionamiento estratégico que puede revertir alineaciones habituales \citep{birch2003tworound}. Cuando el partido del votante queda fuera del balotaje ---como le ocurrió a todos los socios menores de la coalición--- la decisión de cuál de los dos candidatos apoyar queda determinada por factores sustantivos (política pública, liderazgos, identidades) que pueden diferir sistemáticamente entre contextos geográficos.

El presente estudio realiza las siguientes contribuciones:

\begin{enumerate}
    \item \textbf{Documenta la ``rebelión rural''}: los circuitos rurales exhiben tasas de defección de CA al FA 21{,}8 pp superiores a los urbanos ($P(T_{\text{rural}} > T_{\text{urbano}} \mid \text{datos}) > 0{,}99$), el hallazgo más contraintuitivo del estudio.

    \item \textbf{Primera aplicación sistemática de inferencia ecológica bayesiana a nivel circuital en balotajes uruguayos}: 7{,}271 circuitos en 2024 y 7{,}118 en 2019, máxima resolución geográfica disponible en datos oficiales.

    \item \textbf{Identifica la paradoja de la coalición cohesiva derrotada}: pese a mejoras generalizadas en la retención de votos entre 2019 y 2024, la coalición perdió el balotaje ---efecto directo del colapso del caudal de CA ($-$78\%)---, con implicancias para la teoría de coaliciones presidenciales.

    \item \textbf{Valida metodológicamente los resultados}: análisis de sensibilidad con priors Dirichlet genuinamente distintos ($\alpha \in \{0.5, 1, 5\}$) y 500 replicaciones bootstrap confirman la estabilidad de los hallazgos.
\end{enumerate}

El artículo se organiza como sigue. La Sección~\ref{sec:metodologia} describe el modelo bayesiano de transferencia de votos, los datos y la estrategia analítica. La Sección~\ref{sec:resultados} presenta los resultados nacionales y la heterogeneidad geográfica. La Sección~\ref{sec:temporal} examina la evolución 2019--2024. La Sección~\ref{sec:discusion} discute las implicaciones teóricas y las limitaciones del estudio.
